% Options for packages loaded elsewhere
\PassOptionsToPackage{unicode}{hyperref}
\PassOptionsToPackage{hyphens}{url}
\documentclass[
  ignorenonframetext,
]{beamer}
\newif\ifbibliography
\usepackage{pgfpages}
\setbeamertemplate{caption}[numbered]
\setbeamertemplate{caption label separator}{: }
\setbeamercolor{caption name}{fg=normal text.fg}
\beamertemplatenavigationsymbolsempty
% remove section numbering
\setbeamertemplate{part page}{
  \centering
  \begin{beamercolorbox}[sep=16pt,center]{part title}
    \usebeamerfont{part title}\insertpart\par
  \end{beamercolorbox}
}
\setbeamertemplate{section page}{
  \centering
  \begin{beamercolorbox}[sep=12pt,center]{section title}
    \usebeamerfont{section title}\insertsection\par
  \end{beamercolorbox}
}
\setbeamertemplate{subsection page}{
  \centering
  \begin{beamercolorbox}[sep=8pt,center]{subsection title}
    \usebeamerfont{subsection title}\insertsubsection\par
  \end{beamercolorbox}
}
% Prevent slide breaks in the middle of a paragraph
\widowpenalties 1 10000
\raggedbottom
\AtBeginPart{
  \frame{\partpage}
}
\AtBeginSection{
  \ifbibliography
  \else
    \frame{\sectionpage}
  \fi
}
\AtBeginSubsection{
  \frame{\subsectionpage}
}
\usepackage{iftex}
\ifPDFTeX
  \usepackage[T1]{fontenc}
  \usepackage[utf8]{inputenc}
  \usepackage{textcomp} % provide euro and other symbols
\else % if luatex or xetex
  \usepackage{unicode-math} % this also loads fontspec
  \defaultfontfeatures{Scale=MatchLowercase}
  \defaultfontfeatures[\rmfamily]{Ligatures=TeX,Scale=1}
\fi
\usepackage{lmodern}
\ifPDFTeX\else
  % xetex/luatex font selection
\fi
% Use upquote if available, for straight quotes in verbatim environments
\IfFileExists{upquote.sty}{\usepackage{upquote}}{}
\IfFileExists{microtype.sty}{% use microtype if available
  \usepackage[]{microtype}
  \UseMicrotypeSet[protrusion]{basicmath} % disable protrusion for tt fonts
}{}
\makeatletter
\@ifundefined{KOMAClassName}{% if non-KOMA class
  \IfFileExists{parskip.sty}{%
    \usepackage{parskip}
  }{% else
    \setlength{\parindent}{0pt}
    \setlength{\parskip}{6pt plus 2pt minus 1pt}}
}{% if KOMA class
  \KOMAoptions{parskip=half}}
\makeatother
\usepackage{graphicx}
\makeatletter
\newsavebox\pandoc@box
\newcommand*\pandocbounded[1]{% scales image to fit in text height/width
  \sbox\pandoc@box{#1}%
  \Gscale@div\@tempa{\textheight}{\dimexpr\ht\pandoc@box+\dp\pandoc@box\relax}%
  \Gscale@div\@tempb{\linewidth}{\wd\pandoc@box}%
  \ifdim\@tempb\p@<\@tempa\p@\let\@tempa\@tempb\fi% select the smaller of both
  \ifdim\@tempa\p@<\p@\scalebox{\@tempa}{\usebox\pandoc@box}%
  \else\usebox{\pandoc@box}%
  \fi%
}
% Set default figure placement to htbp
\def\fps@figure{htbp}
\makeatother
\setlength{\emergencystretch}{3em} % prevent overfull lines
\providecommand{\tightlist}{%
  \setlength{\itemsep}{0pt}\setlength{\parskip}{0pt}}
\usepackage{bookmark}
\IfFileExists{xurl.sty}{\usepackage{xurl}}{} % add URL line breaks if available
\urlstyle{same}
\hypersetup{
  hidelinks,
  pdfcreator={LaTeX via pandoc}}

\author{\texorpdfstring{}{}}
\date{}

\begin{document}

\begin{frame}{Entropic Architectures: The Historical Instability of Gold
and Silver}
\protect\phantomsection\label{entropic-architectures-the-historical-instability-of-gold-and-silver}
The architecture of British bimetallic entropy was formally codified on
21 September 1717 by the Master of the Royal Mint, Sir Isaac Newton.
Newton, appointed Warden in 1696 through John Locke and Charles Montagu,
and promoted to Master in 1699, approached the currency crisis---the
continuous outflow of English silver to the East India Company and the
continent---with the same geometric precision he applied to the cosmos
in his \emph{Principia}. In his report ``State of the Gold \& Silver
Coins of the Kingdom,'' Newton executed a bureaucratic fiat: he fixed
the golden guinea at exactly 21 silver shillings, establishing a
domestic gold-to-silver ratio of approximately 1:15.21
{[}@quinn1996competing{]}.

This valuation led to a Royal Proclamation on 21 December 1717. Newton
claimed to have brought the coinage to ``a much greater degree of
exactness than ever was known before,'' relying on Trial Plates and
mathematical ratios to formally weigh out the realm
{[}@horsefield1960british{]}. For William Blake, however, this exactness
was not a stabilizing measure but an ontological constraint. Newton's
intervention epitomizes the Urizenic instinct: the attempt to arrest the
living flux of human exchange by superimposing a geometric grid.

\begin{block}{The Alchemy of the Mint and the Geometry of Death}
\protect\phantomsection\label{the-alchemy-of-the-mint-and-the-geometry-of-death}
Newton's role presents a historical paradox that aligns with Blake's
mythological critique. While enforcing the mechanistic ratios of the
1717 standard, Newton was simultaneously immersed in clandestine
alchemical pursuits, experimenting with transmutations and heterodox
theology {[}@schouten2021lucretius{]}. Yet his public legacy at the Mint
stripped currency of its qualitative, alchemical vitality.

Blake isolates this paradox in his 1795 color print \emph{Newton}. The
scientist sits naked at the bottom of the ocean, his gaze fixed downward
upon a geometric scroll, wielding golden compasses---the exact
instruments of the Mint's standardization. As Alan Moore argues in his
2014 Royal Academy commentary on the work, Blake's portrayal is not a
literal portrait but an \emph{``idealized essence''} capturing the
specific financial tyranny Newton exacted as Warden and Master of the
Mint {[}@moore2014blake{]}. Newton had overseen the Great Recoinage of
1696, pursued counterfeiters with unusual severity, and helped entrench
the institutional habits that would later be read as a quiet shift
toward gold dominance.

Just as Urizen creates the material world by division---binding the
infinite with the ``golden compass'' and ``brazen quadrant'' (\emph{The
First Book of Urizen}) {[}@blake1794urizen{]}---Newton's valuation
clamped down upon the living economy. Imposing rigid numerics over the
transformative alchemical process, Newton institutionalized the ``Single
Vision'' Blake fought to break. As Moore characterizes it, Newton's
metrology installed the parameters of an ``internal dungeon'' to which
all humanity was condemned, enthroning a god of mathematical knowledge
whose fiat rendered mankind's living dream-life into pure materialism
{[}@moore2014blake{]}. The print is not merely an aesthetic critique of
empiricism; it is a historically precise indictment of the metrological
regime Newton imposed upon the imperial world, linking mechanical
materialism to the petrification of divine vision
{[}@albernaz2024common; @makdisi2003william{]}.

Furthermore, Newton's 1:15.21 domestic ratio structurally overvalued
gold against fluctuating international bullion markets, which hovered
near 1:14.5 in France and 1:15 across Europe {[}@quinn2007recoinage{]}.
A guinea worth 21 shillings in England was priced at roughly 20s 9d in
Portugal or France. This mathematical misalignment, rather than
fostering stability, guaranteed that gold---particularly Brazilian gold
via Portugal---flooded English ports, while ``undervalued'' silver was
siphoned abroad. By 1715, nearly all the silver Newton struck during the
Great Recoinage (1696--1699) had vanished to France and the East Indies
{[}@quinn1996competing{]}. The Bank of England amassed 15.5 tonnes of
gold by 1730 from these asymmetric flows, inadvertently shifting Britain
toward gold dominance.

\begin{figure}
\centering
\pandocbounded{\includegraphics[keepaspectratio,alt={The Arbitrage Fracture --- Transatlantic Divergence of British, American, and Market Ratios (1717--1873). Three lines trace the divergent monetary architectures: Newton's frozen English Mint Ratio (15.21:1, dashed white), Hamilton's US Mint Ratio (15:1 → 16:1 after 1834, green), and the fluctuating European market reality (gold markers). Fill-between shading encodes the Gresham ``escape velocity''---the precise profit margin that drove bullion arbitrageurs to drain England's circulating silver (silver undervalued, drained abroad) and import gold (gold undervalued, imported). Nine vertical event markers annotate the US Coinage Act (1792), Bank Restriction (1797), British Coinage Act (1816), US ratio revision (1834), Bank Charter Act (1844), and the ``Crime of 1873.'' A shaded band marks the Bank Restriction Period (1797--1821). Sources: Quinn 1996; Flandreau 1996; Laughlin 1898; Horsefield 1960 {[}@flandreau1996bimetallism; @laughlin1898bimetallism{]}.}]{../output/figures/historic_ratio_dark.pdf}}
\caption{The Arbitrage Fracture --- Transatlantic Divergence of British,
American, and Market Ratios (1717--1873). Three lines trace the
divergent monetary architectures: Newton's frozen English Mint Ratio
(15.21:1, dashed white), Hamilton's US Mint Ratio (15:1 → 16:1 after
1834, green), and the fluctuating European market reality (gold
markers). Fill-between shading encodes the Gresham ``escape
velocity''---the precise profit margin that drove bullion arbitrageurs
to drain England's circulating silver (silver undervalued, drained
abroad) and import gold (gold undervalued, imported). Nine vertical
event markers annotate the US Coinage Act (1792), Bank Restriction
(1797), British Coinage Act (1816), US ratio revision (1834), Bank
Charter Act (1844), and the ``Crime of 1873.'' A shaded band marks the
Bank Restriction Period (1797--1821). Sources: Quinn 1996; Flandreau
1996; Laughlin 1898; Horsefield 1960 {[}@flandreau1996bimetallism;
@laughlin1898bimetallism{]}.}\label{fig-historic-ratio}
\end{figure}
\end{block}

\begin{block}{The Triumph of Single Vision in the Bimetallic Standard}
\protect\phantomsection\label{the-triumph-of-single-vision-in-the-bimetallic-standard}
The Newtonian standard did not stabilize value; it maximized the margin
between statutory decree and market reality, fueling the arbitrage it
sought to suppress. It instituted a mechanistic worldview where wealth
derived not from prophetic creation or human labor, but from the
calculation of metallic ratios. This was the triumph of what Blake
termed \emph{Single Vision}---the reduction of qualitative human life to
quantitative, fungible units {[}@erdman1988complete{]}. As he states in
\emph{There is No Natural Religion}, ``He who sees the Ratio only sees
himself only''---diagnosing how this monometallic ratio-setting
collapses perception into narcissistic materialism
{[}@erdman1988complete{]}. Blake famously begged, ``May God us keep /
From Single vision \& Newton's sleep,'' recognizing that such
restrictive geometry severed the living vine of human interaction,
replacing it with the unyielding boundaries of the coin. As Northrop
Frye identifies in \emph{Fearful Symmetry}, Blake's opposition to this
regime was a fully articulated philosophical telos: Fourfold Vision as
the restoration of imaginative creativity operating against the
``tyranny of custom'' and the ``placid ovine herd'' of market-driven
uniformity {[}@frye1947fearful{]}.
\end{block}
\end{frame}

\end{document}
